% -------------------------------------------------------------------
% PAGE-LENGTH KNOBS - tuned to fill about two A4 pages.
%   Ends too early on page 2 -> raise \linespread to 1.1 or 1.15,
%                               or raise the geometry margin to 1.15in.
%   Spills onto page 3       -> lower \linespread to 0.95,
%                               or lower the geometry margin to 0.9in.
% -------------------------------------------------------------------
\documentclass[12pt,a4paper]{article}
\usepackage[T1]{fontenc}
\usepackage[margin=1in]{geometry}
\usepackage{enumitem}
\usepackage{titlesec}
\usepackage[hidelinks]{hyperref}

\linespread{1.0}
\titlespacing*{\section}{0pt}{14pt}{6pt}
\titleformat{\section}{\normalfont\Large\bfseries}{}{0pt}{}
\setlist[itemize]{topsep=5pt,itemsep=4pt,parsep=0pt,leftmargin=*}
\setlist[enumerate]{topsep=5pt,itemsep=4pt,parsep=0pt,leftmargin=*}
\setlength{\parindent}{0pt}
\setlength{\parskip}{8pt}
\pagestyle{empty}

\begin{document}

{\centering\huge\bfseries Silent Tool-Call Failures\par}
\vspace{8pt}
{\centering\Large How Quantization Degrades Function Calling and Structured Output in Small LLMs\par}
\vspace{18pt}

\section*{The Idea}

Local and on-premise agents run quantized models, almost always at 4-bit. Practitioners report a
characteristic failure: conversational quality looks unchanged, standard benchmark accuracy looks
unchanged, and yet a JSON field is silently coerced to the wrong type, an enum value drifts, or a
tool fires with a plausible but incorrect argument. Unlike a badly written sentence, a wrong tool
argument is a wrong database query, a wrong transaction, or a wrong actuator command.

This project measures how much of that degradation is attributable to quantization, separates
failures that are syntactic from failures that are semantic, and tests whether the standard fix,
grammar-constrained decoding, actually addresses the part that matters.

\section*{Datasets and Setup}

\begin{itemize}
  \item \textbf{Berkeley Function Calling Leaderboard (BFCL v4)}, AST categories: simple, multiple,
        parallel, parallel-multiple, plus relevance and irrelevance detection. Open source and
        runnable locally against Hugging Face models.
  \item \textbf{API-Bank} (Level-1 and Level-2 calls) as a second harness, guarding against
        benchmark-specific artefacts.
  \item \textbf{ToolACE}, sampled held-out calls, for a more diverse API surface.
  \item \textbf{A purpose-built schema-stress suite} of roughly 400 cases, each isolating one failure
        class: deeply nested objects, enumerated values, float-versus-integer arguments, optional
        fields, unicode payloads, and menus of fifty candidate tools.
  \item \textbf{Models:} Qwen2.5-1.5B/3B/7B-Instruct, Llama-3.2-1B/3B-Instruct, xLAM-2-1b/8b,
        Hammer2.1, Phi-4-mini.
  \item \textbf{Precisions:} BF16, INT8, NF4, GPTQ-4bit, AWQ-4bit, with INT8/INT4 KV-cache
        quantization as a second axis.
\end{itemize}

\section*{Existing Work}

\begin{itemize}
  \item \textbf{BFCL} (ICML 2025, PMLR 267:48371--48392) is the standard evaluation harness for
        function calling. It compares models, not numerical precisions.
  \item \href{https://arxiv.org/abs/2411.13547}{arXiv:2411.13547} (ToolScan) provides an error
        taxonomy for tool-using models, established at full precision only.
  \item \href{https://arxiv.org/abs/2607.14181}{arXiv:2607.14181} studies quantization for code
        generation and finds that methods differ meaningfully in their effect on correctness. This
        is the nearest analogue in the literature, but the task is different: code is graded by
        execution, tool calls by argument-level correctness.
  \item \href{https://arxiv.org/abs/2505.15030}{arXiv:2505.15030} evaluates on-device LLM
        quantization for accuracy, latency and memory, without examining structured-output validity.
  \item \href{https://arxiv.org/abs/2604.22985}{arXiv:2604.22985} proposes uncertainty quantification
        for function calling, supplying a full-precision baseline verifier to compare against.
  \item Reports of reproducible silent tool-call failures under INT4 KV-cache and W4A16 builds exist
        only as practitioner write-ups. \textbf{No peer-reviewed study measures quantization against
        function calling}, which is an unusual gap for an otherwise crowded area.
\end{itemize}

\section*{Novelty and Contribution}

\begin{enumerate}
  \item The first systematic study of quantization against tool calling, with a failure taxonomy
        that separates \emph{syntactic} breakage (invalid JSON, schema violation) from
        \emph{semantic} breakage (valid schema but wrong tool, wrong argument value, hallucinated
        parameter, or unnecessary tool use on an irrelevant query).
  \item The \emph{constrained-decoding gap}. Grammar-constrained decoding is the accepted remedy for
        malformed output. This work quantifies how much of the quantization penalty it removes,
        with the expectation that syntactic errors disappear while semantic errors persist. That
        single measurement is the paper's headline result.
  \item A scale-asymmetry result. Small models are known to degrade far more than large ones under
        4-bit quantization. Tool calling should amplify this, because argument generation is
        low-entropy and unforgiving, which tells practitioners where the safe parameter floor lies.
  \item Three released artefacts: the schema-stress suite; a cheap call-level verifier combining
        argument-token log-probability margin with a one-shot self-check, targeting detection of
        semantic errors before execution at under twenty percent additional token cost; and a
        reliability-per-gigabyte curve mapping measured memory and latency against call accuracy.
\end{enumerate}

\section*{Real-World Impact}

Agents deployed on-premise or on-device, including privacy-sensitive enterprise assistants, offline
mobile assistants and embedded controllers, all run quantized models with tool access. The engineers
who choose a quantization level for these systems currently rely on chat-quality and general
reasoning benchmarks that do not measure structured-output reliability at all. This work gives them
the missing number, an inexpensive pre-execution check for the failures that remain, and evidence
about which quantization methods are safe for agentic use.

\end{document}
